% !PRINTED 234
\weberpage{234}
%<<CONT

\gesperrt{3. If} $J(x_1, x_2)$ \gesperrt{is any invariant form of the icosahedral group, and} $\xi$ \gesperrt{one of its roots, so that} $J(\xi, 1) = 0$, \gesperrt{then all those quantities are also roots of} $J$ \gesperrt{which arise from} $\xi$ \gesperrt{through the fractional icosahedral substitutions.}

If therefore $J$ has a divisor in common with one of the functions $f$, $H$, $T$, then $J$ is divisible by the form in question.

We begin by factoring the highest possible powers of the three fundamental forms $f$, $H$, $T$ out of $J$, so that $J$ is coprime to these functions. If then $\xi$ is a root of $J$, we can determine $\eta$ so that $\xi$ is a root of the form
\[
\Theta = H^3 - \eta f^5
\]
and $\xi$ is then, by 2., a simple root of $\Theta$. Then, by 3., all the remaining roots of $\Theta$ must also at the same time be roots of $J$, i.e. $J$ is divisible by $\Theta$. The quotient is again an invariant form of the icosahedron, and the argument can be repeated. We thus arrive at the theorem:

\gesperrt{4. Every invariant of the icosahedron can be represented in the form:}
\weq{2}{J(x_1,x_2) = C f^\alpha H^\beta T^\gamma F(H^3, f^5),}
\gesperrt{wherein $\alpha$, $\beta$, $\gamma$ denote non-negative integers, $C$ a constant, and $F$ an integral homogeneous function.}

\begin{center}
\S.\,62.\\
\gesperrt{Polyhedral groups of the second kind. Crystallographic groups.}
\end{center}

We have already, in \S.\,51, called attention to groups of linear ternary substitutions which contain, besides the properly orthogonal substitutions, also improperly orthogonal ones, and which we call \gesperrt{groups of the second kind}.

We now wish to designate the finite ones among them as \gesperrt{polyhedral groups} of the second kind, or also as \gesperrt{extended polyhedral groups}, and the substitutions contained in them by the
%CONT>>
